\section{Introduction}
\label{sec:intro}

% Open model hubs have turned pretrained checkpoints into reusable software artifacts. A developer can now download a model from repositories such as Hugging Face\cite{huggingface} and execute it with only a few lines of code. This convenience also creates a new supply-chain attack surface: a model file may embed hidden behavior triggered during deserialization or inference, even when the surrounding project appears benign\cite{attackblog1}\cite{attackblog2}\cite{attackblog3}\cite{attackblog4}\cite{models_are_codes}\shaofei{need citation, paper/news/blog}. 
% \shaofei{you need some statistics here to show the severity of the problem. For example, what's the impact of recent incidents, how many models are affected, how many users are impacted, etc.}

Malicious code poisoning in the AI model supply chain has emerged as a critical and rapidly growing threat. Unlike traditional software, model files are opaque binary artifacts routinely downloaded and executed with minimal inspection from platforms such as Hugging Face~\cite{huggingface}. Recent attack reports reveal that compromised models have already caused large-scale theft of user credentials and illicit exploitation of LLM API tokens, affecting hundreds of thousands of users and exposing CI/CD secrets of major service providers~\cite{litellm1,litellm2,litellm3}. These incidents demonstrate that a single poisoned model can propagate broadly through the reuse-driven AI ecosystem, amplifying the blast radius of each attack~\cite{models_are_codes}.

% Different from common application code, model files are opaque and generally trusted by users, enabling malicious models to spread widely and cause severe losses. Recent attack reports\cite{litellm1} reveal that such supply chain threats could spread rapidly, affecting hundreds of thousands of end users\cite{litellm2} and posing serious risks of sensitive data leakage, including CI/CD secrets of large model service providers, cloud credentials and user private information\cite{litellm3}.
% \shaofei{This paragraph introduce the model supply chain attack(maybe need a professtional term). How do you classify it? Introduce each type of attack and find references. }
% To systematically understand these AI supply-chain vulnerabilities, existing threats can be broadly classified into two categories of malicious code poisoning attacks. The first and historically most prevalent is insecure serialization attacks. Major deep learning frameworks frequently rely on expressive serialization formats to save complex model states. Attackers exploit this by embedding malicious payloads that execute arbitrary code automatically during the deserialization process\cite{liu2025arthideseekmaking}. Recent studies show that insecure serialization remains widespread on public model hubs, and that malicious model artifacts can be published and propagated through ordinary sharing workflows \cite{3765037,casey2024largescaleexploitinstrumentationstudy}. The second category is framework-native attacks. Rather than exploiting file format vulnerabilities, attackers weaponize legitimate, built-in APIs of the AI frameworks themselves to perform malicious operations. For example, attackers maliciously design the model's computational graph and embed dangerous framework operators directly into the inference pipeline\cite{zhu2025}. Consequently, these stealthy payloads execute under the disguise of normal tensor operations and can effectively evade traditional static security scanners.

To defend against these threats, existing research has focused on static detection methods that identify risks by scanning serialized files for known malicious signatures, decompiling pickle files, or detecting suspicious imports~\cite{modelscan,fickling,picklescanning}. These tools reason about code structure rather than executed behavior, frequently confusing benign model operations with malicious attacks. For instance, if a benign model fetches remote weights and a malicious payload creates a reverse shell using the same networking library, such tools cannot distinguish the two~\cite{nabeel2026deepdiveabusedl}.
Therefore there is a growing consensus that dynamic detection is necessary~\cite{10068497,10100797898195307244}. However, existing dynamic methods rely on prior knowledge of attack patterns to instrument specific functions or filter system events, making them vulnerable to rapidly evolving adversarial techniques that embed payloads into computational graphs using only legitimate framework APIs~\cite{zhu2025}. Static tools scan for known signatures or decompile bytecodes; dynamic methods hard-code which functions to hook based on expert heuristics and depend on static whitelists~\cite{299852}. Because adversarial techniques iterate rapidly, any defense contingent on knowing what to look for is inherently vulnerable, suffering from high false-positive and false-negative rates and complete defenselessness against zero-day exploits~\cite{GUO2023175, aghakhani}.



% Existing defenses against such threats share a fundamental and underexamined limitation: they require prior knowledge of attack patterns to function effectively. Static detection tools scan serialized files for known malicious signatures, import dangerous-function blacklists, or decompile pickle bytecodes to identify suspicious patterns~\cite{modelscan,fickling,picklescanning}. These tools reason about code structure rather than executed behavior, so they cannot distinguish a benign model loading remote weights from a malicious payload establishing a reverse shell when both invoke the same networking library. The problem is even more severe for dynamic methods: instrumentation techniques hard-code which specific functions to hook based on expert heuristics, and system-call monitors depend on static whitelists to filter benign events~\cite{299852,294545}. Because adversarial techniques iterate rapidly—attackers now embed payloads directly into computational graphs using only legitimate framework APIs~\cite{zhu2025}—any defense contingent on knowing what to look for is inherently vulnerable. Such systems not only produce high false-positive and false-negative rates, but are fundamentally defenseless against zero-day exploits that lack established signatures.


% % Unlike ordinary application code, model artifacts are opaque, framework-dependent, and routinely trusted by downstream users, which makes pre-deployment inspection particularly difficult. 
% However, the majority of existing AI security research has predominantly focused on content safety and robustness, such as data poisoning and jailbreak attacks~\cite{3708531,wang2025sokunderstandingvulnerabilitieslarge, 11029806, 11029795}. Consequently, a significant research gap remains in detecting system-level threats within the AI model supply chain\cite{3690338}\cite{digregorio2026insecurityloadingmachinelearning}.
% Existing solutions against such threats primarily adopt static detection methods, which identify potential risks by scanning serialized files, decompiling pickle files, or detecting suspicious import behaviors\cite{modelscan}\cite{fickling}\cite{picklescanning}. These tools are useful as a first barrier, but they reason about code structure rather than executed behavior. As a result, they frequently generate false positives against benign programs while still missing attacks hidden behind legitimate APIs or code obfuscation. This limitation becomes particularly severe for camouflage attacks such as TensorAbuse\cite{zhu2025}, where an attacker use authorized TensorFlow interfaces to perform file access, network exfiltration, code injection, or shell acquisition without obviously malicious syntax.
% % Therefore, it is more effective to focus on whether the model performs malicious operations at runtime. However, a runtime defense for model artifacts is nontrivial. First, source code instrumentation relies on the availability of the model's source code, which is often proprietary or inaccessible. Moreover, while it captures dynamic behaviors at the application layer, it fails to penetrate the underlying library APIs, rendering it completely blind to sophisticated attacks embedded within the low-level implementations. Second, binary instrumentation incurs prohibitive performance overhead and high deployment costs. Given the vast and highly fragmented ecosystem of Python dependencies used in modern AI frameworks, it is practically impossible to pre-instrument all underlying libraries to construct a comprehensive and automated detection system. Third, system call monitoring suffers from a severe semantic gap. Low-level system events lack the necessary high-level context to directly deduce the model's actual intent, inevitably leading to massive false positives as benign operations like logging and dependency loading are easily confused with genuine exfiltration and persistence.
% Therefore, it is more effective to focus on whether the model performs malicious operations at runtime. However, establishing a robust runtime defense for model artifacts remains nontrivial. A fundamental limitation of existing dynamic detection methods is their strict reliance on priori knowledge\cite{299852}\cite{294545}. For instance, current instrumentation techniques typically require predefined attack signatures or expert-guided heuristics to determine which specific functions to hook. Similarly, system call monitors heavily depend on static whitelists to filter out benign events. Because adversarial techniques iterate rapidly, such defenses relying on prior knowledge possess inherent vulnerability. They not only suffer from high rates of false positives and negatives, but are also fundamentally defenseless against zero-day exploits that lack established signatures.

To overcome this reliance on prior knowledge, we shift our strategy from hunting known threats to modeling benign behavioral characteristics. Our key observation is that benign AI models exhibit highly regular runtime templates, governed by the framework's internal logic and largely invariant across architectures and tasks. Normal execution consistently passes through a small set of stable phases such as deserialization, tensor materialization, and library loading. Malicious models inevitably deviate from these templates, producing anomalous patterns such as unexpected file access during deserialization or unauthorized network communication during inference.

Building on this observation, we propose \toolname{}, a defense that requires no prior knowledge of specific attack methodologies. \toolname{} transforms dynamic model security protection into a two-stage behavior analysis problem, bridging the semantic gap between low-level system evidence and model intent. In the first stage, an unsupervised model learns the distribution of benign operation patterns from system call traces and filters out normal activity, requiring no attack signatures or whitelists. Only statistically abnormal operations are escalated to the second stage, where they are associated with Python semantic context collected via lightweight, non-invasive instrumentation. This cross-layer association combines kernel-level system call evidence with application-layer semantics. Finally, an LLM traces the root cause of underlying operations and judges whether they conform to the regular execution logic of the model or constitute malicious attacks. Through this two-stage pipeline, \toolname{} achieves both high detection accuracy and low false-positive rates while maintaining practical efficiency.


However, realizing \toolname{} faces three non-trivial technical challenges. First, how to define a meaningful unit of system call behavior that can be learned by a machine learning model. A single system call is excessively fine-grained and massive in volume, while malicious behaviors generally manifest as consecutive correlated call sequences\cite{3677893,jcp6030099}. \toolname{} addresses this by aggregating system calls into lifecycle-bounded operations, converting raw traces into semantically meaningful operations.
Second, how to achieve comprehensive semantic instrumentation across heterogeneous, rapidly evolving AI frameworks without source-code access or prohibitive overhead. \toolname{} leverages an LLM to dynamically generate framework-adaptive monkey-patching scripts that automatically discover and wrap public callable functions at runtime.
Third, how to reliably align Python-level semantic events with kernel-level system call traces under high concurrency, where naive timestamp matching produces spurious associations. \toolname{} addresses this through a fuzzy spatio-temporal alignment strategy that combines temporal windowing with resource-related spatial hints to robustly associate each abnormal operation with its true application-layer context.

% \toolname{} employs an unsupervised Autoencoder to statistically screen for deviations from learned benign operation distributions, filtering out normal activity without relying on attack signatures or whitelists. Only statistically abnormal operations are escalated to a cross-layer analysis stage, where they are associated with high-dimensional Python semantic context collected via lightweight, non-invasive instrumentation. This association bridges the semantic gap between low-level system events and model intent, enabling LLM-based reasoning to determine whether observed behavior is benign or indicative of an attack. Through this two-stage pipeline, \toolname{} achieves high detection accuracy and low false-positive rates while maintaining practical efficiency.



% ===== ORIGINAL DRAFT PARAGRAPHS (comment preserved) =====
% To address the inherent limitations of such defense strategies, we shift our focus from hunting known threats to modeling benign behavioral characteristics. Our main observation is that benign AI models exhibit highly regular runtime templates. These execution phases are dictated by the framework's internal logic and remain largely invariant regardless of the model's specific architecture or task. Normal execution consistently passes through a small set of stable phases, including deserialization, tensor materialization, and library loading. In contrast, malicious models inevitably deviate from these strict templates, exhibiting anomalous patterns such as unexpected file access during deserialization or unauthorized network communication during inference. 
% Building on this observation, we propose a defense mechanism that requires no prior knowledge.
% We first employ an unsupervised, statistical approach to screen for deviations from these regular benign templates. To capture the necessary context without imposing filtering biases, we implement a lightweight, universal instrumentation strategy that hooks all library API calls. We completely abandon unstable whitelist filtering to avoid poor performance caused by reliance on prior knowledge. By combining statistical baseline profiling with comprehensive and low-overhead monitoring, our approach effectively defend against both known and zero-day threats without requiring any prior knowledge of specific attack methodologies.
% Specifically, \toolname{} leverages this regularity to transform dynamic model security protection into a two-stage behavior analysis problem. In the first stage, an unsupervised anomaly detector learns the distribution of benign operation patterns and efficiently filters out normal activity. In the second stage, only statistically abnormal operations are submitted for cross-layer analysis, where they are associated with high-level Python semantic context collected by adaptive instrumentation. This association bridges the semantic gap between low-level system events and model intent, enabling LLM-based reasoning to determine whether the observed behavior is consistent with benign execution or indicative of an attack. Through this two-stage pipeline, \toolname{} achieves both high detection accuracy and low false positive rates while maintaining practical efficiency.

Experiments on 89 models from the public dataset and 120 models from the advanced synthetic dataset demonstrate that \toolname{} achieves 100\% precision, 98.70\% F1-score, and 0\% false positives on benign models using sensitive framework APIs, while introducing only 5.69\% average runtime overhead. In summary, our contributions are:
  
\begin{itemize}
\item We propose the first hybrid dynamic detection system for AI model supply chain attacks that combines non-bypassable system-call evidence with Python semantic context, bridging the semantic gap between low-level behaviors and model intent.
\item We develop a robust cross-layer association mechanism that aligns abnormal system-level operations with high-level semantic events, enabling LLM-based reasoning to provide accurate and explainable security alerts.
\item Extensive experiments on real-world and synthetic malicious models demonstrate that our system achieves high detection accuracy and low false-positive rates, outperforming existing static defenses with good practical deployment value.
\item We open source our system and provide a public benchmark of malicious models at \url{https://anonymous.4open.science/r/DDMO-F184/} to facilitate future research in AI supply chain security.
\end{itemize}

